\chapter*{Conclusion}
\addcontentsline{toc}{chapter}{Conclusion}
\label{ch:conclusion}

{}

The study have assessed the impacts of urbanization on plant diversity and ecosystem functioning across urbanization gradientsin Makurdi and Otukpo Local Government Areas of Benue State, Nigeria. By integrating land-use/land-cover analysis, field-based vegetation surveys, soil assessments, and statistical modelling, the study provides a comprehensive understanding of how rapid urban expansion alters ecological patterns in tropical urban environments.

The results revealed that built-up areas increased by over 517\% between 1990 and 2024, replacing extensive forest and cropland. This large scale conversion have also led to the following; the habitat fragmentation, reduced ecological complexity, and altered hydrological patterns. Plant diversity exhibited a clear decreasing trend along the urban gradient, with rural areas maintaining the highest richness and Shannon diversity, while urban zones showed simplified plant communities dominated by disturbance-tolerant species. Also, the soil properties as revealed by the findings of the study, showed the urban-induced modifications, particularly a significant increase in soil pH, indicating disruption of natural nutrient processes. Vegetation indices (NDVI/EVI) further confirmed substantial declines in vegetation health and primary productivity toward urban Centre's. The regression modelling in its findings have demonstrated that population density was the strongest suppressor of vegetation greenness, highlighting demographic pressure as a major ecological driver. These findings collectively demonstrate that urbanization in Benue State is driving a transition from biologically rich rural ecosystems to ecologically degraded urban environments. The ecological consequences include biodiversity loss, reduced vegetation functioning, altered soil processes, and weakened ecosystem resilience.

Despite its strengths, this study has several limitations that should be acknowledged. First, is the temporal resolution of the remote sense data which is constrained by the availability of cloud-free Landsat imagery for some years, which may have limited the precision of finer-scale temporal analyses. Second, is the plant diversity data and assessment that was conducted only in selected plots and did not incorporate seasonal variations, potentially influencing the detection of species that appear or decline at different times of the year. Third, is the soil data and the analyses which focused primarily on key physicochemical indicators and did not include microbiological assessments part of a whole, which are critical for fully understanding soil ecological functioning and nutrient cycling processes. Additionally, some urbanization indicators such as built-up intensity and night-time light data may not adequately capture the complexity of informal settlement patterns common in African cities. Finally, the regression modelling employed ward-level data aggregation, which, although useful for broad analysis, may obscure important micro-spatial variations occurring within individual neighborhoods or land parcels.

To deepen the understanding of urban ecological dynamics in tropical cities, future research should adopt several methodological improvements. First, the high-resolution satellite datasets like the Sentinel-2 and the WorldView should be incorporated and this is to enhance the precision of land-use/land-cover change detection. Multi-seasonal vegetation sampling is needed to capture temporal fluctuations in species composition and plant community structure. Secondly, the soil aspect of the study should be expanded to include micro-biome and enzyme activity analyses where lab test will carried out, and will provide more detailed insights into ecosystem processes and soil health. Future work should also employ landscape fragmentation metrics such as edge density and patch cohesion to quantify ecological connectivity and habitat integrity. And lastly, the urban ecological modelling tools like InVEST and the landscape metrics should be used to predict potential ecological outcomes under varying urban growth scenarios. In conclusion, the investigation of the socio-economic causes behind the loss of vegetation and that of the land conversion would be providing a more holistic and realistic understanding of urbanization processes and support more integrated and effective planning strategies.

